\documentclass[11pt, aspectratio=169]{beamer}

% LuaLaTeXで日本語を使うための必須パッケージ
\usepackage{luatexja}

% 便利な追加パッケージ
\usepackage{booktabs}   % 表の罫線
\usepackage{tikz}       % 図形描画
\usepackage{listings}   % ソースコード挿入
\lstset{
  basicstyle=\ttfamily\scriptsize,
  frame=single,
  breaklines=true,
  numbers=left,
  numberstyle=\tiny
}

% テーマとカラーテーマの設定
\usetheme{Madrid}
\usecolortheme{orchid}

% ==========================================
% 【改良】セクション・サブセクションの可視化設定
% ==========================================

% 1. セクション開始時の自動表紙（大きな扉絵）
\AtBeginSection[]{
  \begin{frame}
    \vfill
    \centering
    \begin{beamercolorbox}[sep=8pt,center,shadow=true,rounded=true]{title}
      \usebeamerfont{title}\insertsectionhead\par%
    \end{beamercolorbox}
    \vfill
  \end{frame}
}

% 2. サブセクション開始時に「現在のセクション内の目次」を自動挿入
% 今いるサブセクション以外を半透明（shaded）にすることで、現在地が強調されます
\AtBeginSubsection[]{
  \begin{frame}{\insertsectionhead}
    \tableofcontents[currentsection, currentsubsection, subsectionstyle=show/shaded/hide]
  \end{frame}
}

% タイトル情報
\title{妨害通信の全体像と今後の予定}
\subtitle{妨害通信}
\author{学籍番号(Student ID): 1280391 \\ 氏名(Name): 細川夏風(Hosokawa Natsuka)}
\institute{高知工科大学(Kochi Univ of Tech)}
\date{\today}

\begin{document}

\usetikzlibrary{decorations.pathmorphing}

% タイトルスライド
\begin{frame}
  \titlepage
\end{frame}

\begin{frame}{進捗}
  \begin{enumerate}[1]
    \item SAGINの通信妨害
    \item SAGINの概要
    \item 通信妨害と概要(今回)
  \end{enumerate}
\end{frame}

% 全体目次スライド
\begin{frame}{目次}
  \tableofcontents
\end{frame}

% ==========================================
% ここから本文
% ==========================================

\section{通信妨害とは}
\subsection{概要}

% 以下の2行をプリアンブル（\begin{document}の前）に追加してください
% \usepackage{tikz}
% \usetikzlibrary{decorations.pathmorphing}

\begin{frame}{概要}
  \begin{columns}[c, onlytextwidth]
    
    % --- 左側：テキスト部 ---
    \begin{column}{0.52\textwidth}
      無線通信では共通の空間を利用し通信しているため妨害が容易で，守ることも難しい．
    
      \vspace{1em}
    
      ノイズが大きすぎると受信機は$1$, $0$の判別ができなくなり，通信が行えない．
    
      \vspace{1.5em}
    
      \begin{block}{要約}
        妨害通信とは，みんなが共有している空間に悪意ある大量の妨害電波を送信し，正規の電波をかき消す攻撃のこと．
      \end{block}
    \end{column}
    
    \hfill
    
    % --- 右側：TikZ図解部 ---
    \begin{column}{0.5\textwidth}
      \centering
      \begin{tikzpicture}[scale=0.85, transform shape]
        % 共通空間の枠
        \draw[dashed, gray!60, rounded corners, fill=gray!2] (-0.6, -1.3) rectangle (3.8, 2.5);
        \node at (1.6, -1.0) [gray, font=\tiny] {共通の空間（空気中）};

        % 1. 正規送信機（左）
        \draw[thick, blue!80!black] (0,-0.5) -- (0,0.3);
        \draw[thick, blue!80!black, fill=blue!10] (0,0.3) -- (-0.2,0.6) -- (0.2,0.6) -- cycle;
        \node at (0,-0.8) [blue!80!black, font=\scriptsize, align=center] {正規\\送信機};

        % 2. 受信機（右）
        \draw[thick, darkgray] (3.2,-0.5) -- (3.2,0.3);
        \draw[thick, darkgray, fill=gray!20] (3.2,0.3) -- (3.0,0.6) -- (3.4,0.6) -- cycle;
        \node at (3.2,-0.8) [darkgray, font=\scriptsize, align=center] {受信機};

        % 3. 妨害機（上）
        \draw[thick, red!80!black] (1.6,2.0) -- (1.6,1.4);
        \draw[thick, red!80!black, fill=red!10] (1.6,1.4) -- (1.4,1.1) -- (1.8,1.1) -- cycle;
        \node at (1.6,2.3) [red!80!black, font=\scriptsize, align=center] {妨害機\\（悪意ある送信）};

        % 4. 正規の電波（きれいなサイン波）
        % x=0.2から2.8の間で綺麗に4周期波を打つように調整
        \draw[blue, thick, domain=0.2:2.8, samples=100] plot (\x, {0.2*sin((\x-0.2)/2.6 * 4 * 360)});
        \draw[blue, thick, ->] (2.8,0) -- (3.1,0);
        \node at (1.2, 0.4) [blue, font=\tiny] {正規の電波};

        % 5. 妨害電波（ギザギザのノイズ波）
        \draw[red, thick, ->, decorate, decoration={zigzag, amplitude=1.5mm, segment length=2mm}] (1.6, 1.0) -- (3.0, 0.2);
        \node at (2.5, 1.0) [red, font=\tiny, align=center, rotate=-30] {大量の\\妨害電波};

        % 6. 受信機での混信状態のラベル
        \node at (3.2, 1.2) [red!90!black, font=\bfseries\tiny, align=center] {混信ノイズ！\\$1$, $0$ の\\判別不能};
      \end{tikzpicture}
    \end{column}
    
  \end{columns}
\end{frame}

\subsection{通信妨害の種類}

\begin{frame}{通信妨害の種類}
  \begin{columns}[c, onlytextwidth]
    
    % --- 左側：テキスト部 ---
    \begin{column}{0.5\textwidth}
      \begin{itemize}
        \item \textbf{継続的妨害}:\\
          常に強い妨害電波を出し続けるもの．

          \vspace{1em}

          \alert{強力だが，攻撃者のデバイスのバッテリーも消費が激しい．}
      \end{itemize}
    \end{column}
    
    \hfill
    
    % --- 右側：TikZ図解部 ---
    \begin{column}{0.45\textwidth}
      \centering
      \begin{tikzpicture}[scale=0.9, transform shape]
        % 妨害機
        \node[draw=red!80!black, thick, fill=red!10, rounded corners, align=center, inner sep=5pt] (attacker) at (0, 2) {妨害機\\(攻撃者)};
        
        % 受信機
        \node[draw=darkgray, thick, fill=gray!10, rounded corners, align=center, inner sep=5pt] (receiver) at (3.5, 0) {受信機};
        
        % バッテリー
        \node[draw=black, thick, fill=white, align=center, font=\scriptsize] (battery) at (0, 0) {バッテリー\\(消費大)};
        
        % バッテリーから妨害機への線
        \draw[->, thick, gray, dashed] (battery.north) -- (attacker.south);
        
        % 妨害電波の矢印
        \draw[->, red, thick, decorate, decoration={zigzag, amplitude=1.5mm, segment length=2mm}] 
          (attacker.east) to[out=0, in=120] node[above, sloped, font=\scriptsize, text=black] {絶え間ない送信} (receiver.north west);
      \end{tikzpicture}
    \end{column}
    
  \end{columns}
\end{frame}

\begin{frame}{通信妨害の種類}
  \begin{columns}[c, onlytextwidth]
    
    % --- 左側：テキスト部 ---
    \begin{column}{0.52\textwidth}
      \begin{itemize}
        \item \textbf{反応的妨害}: \\
          普段は何も行わないが，通信を開始した瞬間に妨害電波を用いて妨害を行う．

        \vspace{1em}

        \alert{エネルギー効率が良く，妨害時以外は電波を出さないため，発見が困難．}
      \end{itemize}
    \end{column}
    
    \hfill
    
    % --- 右側：TikZ図解部 ---
    \begin{column}{0.45\textwidth}
      \centering
      \begin{tikzpicture}[scale=0.9, transform shape]
        % 正規送信機
        \node[draw=blue!80!black, thick, fill=blue!10, rounded corners, align=center, inner sep=5pt] (sender) at (0, 0) {正規\\送信機};
        
        % 受信機
        \node[draw=darkgray, thick, fill=gray!10, rounded corners, align=center, inner sep=5pt] (receiver) at (3.5, 0) {受信機};
        
        % 妨害機
        \node[draw=red!80!black, thick, fill=red!10, rounded corners, align=center, inner sep=5pt] (attacker) at (1.75, 2.5) {妨害機\\(待機/省電力)};
        
        % 1. 正規の通信
        \draw[->, thick, blue] (sender.east) -- node[below, font=\scriptsize] {① 通信開始} (receiver.west);
        
        % 2. 通信の検知
        \draw[->, thick, gray, dashed] (sender.north) to[out=90, in=180] node[above left, font=\scriptsize, text=black] {② 検知} (attacker.west);
        
        % 3. 妨害電波
        \draw[->, red, thick, decorate, decoration={zigzag, amplitude=1.5mm, segment length=2mm}] 
          (attacker.east) to[out=0, in=90] node[above right, font=\scriptsize, text=black, align=center] {③ 即座に\\妨害！} (receiver.north);
      \end{tikzpicture}
    \end{column}
  \end{columns}
\end{frame}

\begin{frame}{通信妨害の種類}
  \begin{columns}[c] % [c]は左右のコンテンツを上下中央揃えにするオプションです
    
    % --- 左側：テキスト部分（幅55%） ---
    \begin{column}{0.55\textwidth}
      \begin{itemize}
        \item \textbf{スマートジャミング}: \\
          通信データのすべてに妨害をかけるのでなく，通信確立のための「同期信号」，「制御信号」を狙って行う妨害．
          
          \vspace{1em}
          \alert{\texttt{Wi-Fi}や$5$\texttt{G}に用いられる．}
      \end{itemize}
    \end{column}
    
    % --- 右側：図の部分（幅45%） ---
    \begin{column}{0.45\textwidth}
      \centering
      \begin{tikzpicture}[scale=0.9, transform shape]
        % パケット全体の描画
        % 1. 同期/制御信号のブロック（攻撃対象）
        \draw[thick, fill=blue!20] (0, 0) rectangle (1.5, 0.8);
        \node at (0.75, 0.4) [font=\scriptsize, align=center] {同期/制御\\信号};
        
        % 2. データ本体のブロック
        \draw[thick, fill=gray!20] (1.5, 0) rectangle (4.5, 0.8);
        \node at (3.0, 0.4) [font=\scriptsize] {データ本体};
        
        % 妨害機
        \node[draw=red!80!black, thick, fill=red!10, rounded corners, align=center, inner sep=5pt] (attacker) at (0.75, 2.8) {妨害機\\(ピンポイント)};
        
        % 妨害電波（制御信号のみを狙う）
        \draw[->, red, thick, decorate, decoration={zigzag, amplitude=1.5mm, segment length=2mm}] 
          (attacker.south) -- node[left, font=\tiny, text=black, align=right] {ここだけ\\狙い撃ち} (0.75, 0.8);
        
        % データはスルーしていることを示す線（オプション）
        \draw[->, darkgray, dashed] (attacker.east) to[out=0, in=90] node[right, font=\tiny, text=darkgray] {データは無視} (3.0, 0.8);
        
        % 破壊エフェクト（バツ印）
        \draw[ultra thick, red] (0.2, 0.1) -- (1.3, 0.7);
        \draw[ultra thick, red] (0.2, 0.7) -- (1.3, 0.1);
        
        % 結果の注釈
        \node at (2.25, -0.5) [font=\scriptsize\bfseries, red!80!black] {制御不能 $\rightarrow$ 通信全体が成立しない};
      \end{tikzpicture}
    \end{column}
    
  \end{columns}
\end{frame}

\subsection{まとめ}

\begin{frame}{まとめ}
  \begin{block}{}
    これまでに記したように，通信妨害には様々な種類があり，それぞれに一長一短と利用に適した場面がある．
  \end{block}
\end{frame}

\section{通信妨害の対策}

\subsection{対策として用いられている手法}

\begin{frame}{対策手法}
  \begin{columns}[c]
    
    % --- 左側：テキスト部分（幅55%） ---
    \begin{column}{0.55\textwidth}
      \begin{itemize}
        \item \textbf{周波数ホッピング}: \\
          通信に使う周波数をあらかじめ決めた順番で高速に切り替えながら通信を行う．
        
          \vspace{1.5em}
    
        \alert{攻撃者が現在通信している周波数帯を割り出すことが困難．\texttt{Bluetooth}などで用いられている．}
      \end{itemize}
    \end{column}
    
    % --- 右側：図の部分（幅45%） ---
    \begin{column}{0.45\textwidth}
      \centering
      \begin{tikzpicture}[scale=0.9, transform shape]
        % 軸の描画
        \draw[->, thick, darkgray] (0,0) -- (4.5,0) node[right, font=\scriptsize] {時間};
        \draw[->, thick, darkgray] (0,0) -- (0,3.5) node[above, font=\scriptsize] {周波数};
        
        % チャンネルの補助線
        \foreach \y in {1,2,3,4} {
          \draw[dashed, gray!40] (0,\y*0.7) -- (4.2,\y*0.7);
          \node[left, font=\tiny, text=gray] at (0,\y*0.7) {ch\y};
        }
        
        % 攻撃者による妨害電波（特定の周波数を狙い続けると仮定）
        \fill[red!15] (0, 1.4) rectangle (4.2, 2.1);
        \node[red!80!black, font=\tiny, align=left] at (4.0, 2.3) {妨害電波};
        
        % ホッピングする通信データ（青いブロック）
        % t1, ch4
        \fill[blue!60, rounded corners=1pt] (0.4, 2.8) rectangle (0.8, 3.2); 
        % t2, ch1
        \fill[blue!60, rounded corners=1pt] (1.2, 0.7) rectangle (1.6, 1.1); 
        % t3, ch2
        \fill[blue!60, rounded corners=1pt] (2.0, 1.4) rectangle (2.4, 1.8); 
        
        % 妨害と衝突するデータ（ch3のタイミングで妨害に当たる）
        \fill[blue!60, rounded corners=1pt] (2.8, 1.75) rectangle (3.2, 2.15); 
        \draw[thick, red] (2.8, 1.75) -- (3.2, 2.15); % バツ印（通信失敗）
        \draw[thick, red] (2.8, 2.15) -- (3.2, 1.75);
        
        % t5, ch1
        \fill[blue!60, rounded corners=1pt] (3.6, 0.7) rectangle (4.0, 1.1); 
        
        % ホッピングの軌跡（矢印）
        \draw[->, thick, darkgray, dashed] (0.8, 3.0) -- (1.2, 0.9);
        \draw[->, thick, darkgray, dashed] (1.6, 0.9) -- (2.0, 1.6);
        \draw[->, thick, darkgray, dashed] (2.4, 1.6) -- (2.8, 1.95);
        \draw[->, thick, darkgray, dashed] (3.2, 1.95) -- (3.6, 0.9);
        
        % 注釈
        \node[font=\scriptsize\bfseries, blue!80!black] at (2.2, -0.6) {高速に切り替えて通信を確保};
        \node[font=\tiny, red!80!black, align=center] at (3.0, 1.4) {一部だけ\\損失};
      \end{tikzpicture}
    \end{column}
  \end{columns}
\end{frame}

\begin{frame}{対策手法}
  \begin{columns}[c]
    
    % --- 左側：テキスト部分（幅53%） ---
    \begin{column}{0.53\textwidth}
      \begin{itemize}
        \item \textbf{スペクトラム拡散}: \\
          送りたいデータを本来必要な帯域よりも更に広い周波数帯に引き伸ばして送信する方法．
        
          \vspace{1.5em}
    
        \alert{妨害者が強いノイズを出しても受信側でデータを圧縮して戻す際に妨害ノイズが薄まって消える．ノイズに強い．}
      \end{itemize}
    \end{column}
    
    % --- 右側：図の部分（幅47%） ---
    \begin{column}{0.47\textwidth}
      \centering
      \begin{tikzpicture}[scale=0.85, transform shape]
        % 1. 送信時のデータ（狭帯域・高電力）
        \draw[->, thick, darkgray] (0, 0) -- (3.5, 0) node[right, font=\tiny] {周波数};
        \draw[->, thick, darkgray] (0, 0) -- (0, 1.5) node[above, font=\tiny] {電力};
        \fill[blue!70] (1.3, 0) rectangle (1.7, 1.2);
        \node[font=\scriptsize, align=center] at (1.5, 1.5) {元のデータ\\(送信側)};
        
        % 拡散の矢印
        \draw[->, thick, darkgray, dashed] (1.5, -0.2) -- (1.5, -0.8) node[midway, right, font=\scriptsize] {拡散して送信};
        
        % 2. 空間での状態（広い帯域・低電力 ＋ 強力なノイズ）
        \begin{scope}[yshift=-2.6cm]
          \draw[->, thick, darkgray] (0, 0) -- (3.5, 0) node[right, font=\tiny] {周波数};
          \draw[->, thick, darkgray] (0, 0) -- (0, 1.5);
          % 拡散された信号（広く浅く）
          \fill[blue!30] (0.2, 0) rectangle (2.8, 0.3);
          \node[font=\tiny, text=blue!80!black] at (0.8, 0.4) {拡散信号};
          % 妨害ノイズ（特定の周波数に強く出る）
          \fill[red!70] (2.0, 0) rectangle (2.4, 1.3);
          \node[font=\tiny, text=red!80!black] at (2.2, 1.5) {妨害ノイズ};
        \end{scope}
        
        % 逆拡散の矢印
        \draw[->, thick, darkgray, dashed] (1.5, -2.8) -- (1.5, -3.4) node[midway, right, font=\scriptsize] {受信して逆拡散};
        
        % 3. 受信時の状態（データ復元 ＋ ノイズ拡散）
        \begin{scope}[yshift=-5.2cm]
          \draw[->, thick, darkgray] (0, 0) -- (3.5, 0) node[right, font=\tiny] {周波数};
          \draw[->, thick, darkgray] (0, 0) -- (0, 1.5);
          % 復元された信号（元の狭帯域・高電力に戻る）
          \fill[blue!70] (1.3, 0) rectangle (1.7, 1.2);
          \node[font=\tiny, text=blue!80!black] at (1.5, 1.4) {データ復元};
          % 拡散されたノイズ（広く浅く分散される）
          \fill[red!30] (0.2, 0) rectangle (2.8, 0.2);
          \node[font=\tiny, text=red!80!black, align=center] at (2.6, 0.4) {ノイズは\\薄まる};
        \end{scope}
      \end{tikzpicture}
    \end{column}
  \end{columns}
\end{frame}

\begin{frame}{対策手法}
  \begin{columns}[c]
    
    % --- 左側：テキスト部分（幅55%） ---
    \begin{column}{0.55\textwidth}
      \begin{itemize}
        \item \textbf{\texttt{MIMO}と空間フィルタリング}: \\
          スマホや基地局にある複数のアンテナを用いる方法．
    
        \vspace{1em}
    
        複数のアンテナを用いると\alert{正規の信号が来ている方向と妨害電波が来ている方向を計算できる}．そして，受信機で妨害電波が来ている方向からの電波を遮断すればいい．
      \end{itemize}
    \end{column}
    
    % --- 右側：図の部分（幅45%） ---
    \begin{column}{0.45\textwidth}
      \centering
      \begin{tikzpicture}[scale=0.85, transform shape]
        % 受信エリア（正規の方向へ向けるビーム）
        \fill[blue!10] (2.3, 0.4) -- (-0.2, 2.5) arc (130:170:2.5) -- cycle;
        \node[font=\tiny, text=blue!80!black] at (0.6, 1.2) {受信エリア};

        % 受信機（複数アンテナ）
        \node[draw=darkgray, thick, fill=gray!10, rounded corners, minimum width=1.6cm, minimum height=0.6cm, align=center, font=\scriptsize] (rx) at (2.5, 0) {受信機};
        
        % アンテナ素子
        \foreach \x in {1.9, 2.3, 2.7, 3.1} {
            \draw[thick, darkgray] (\x, 0.3) -- (\x, 0.8);
            \fill[darkgray] (\x, 0.8) circle (1.5pt);
        }
        \node[font=\tiny] at (2.5, 1.1) {複数アンテナ};

        % 1. 正規の通信（上から）
        \node[draw=blue!80!black, thick, fill=blue!10, rounded corners, inner sep=3pt, font=\scriptsize] (tx) at (-0.5, 2.2) {正規の電波};
        \draw[->, thick, blue] (tx.east) -- node[above, sloped, font=\tiny, yshift=2pt] {方向を特定} (1.6, 0.9);

        % 2. 妨害機（下から）
        \node[draw=red!80!black, thick, fill=red!10, rounded corners, inner sep=3pt, font=\scriptsize] (jammer) at (-0.5, -1.6) {妨害機};
        
        % ★修正点1：yshift=-8pt を追加し、ギザギザの線から文字を遠ざけました
        \draw[->, thick, red, decorate, decoration={zigzag, amplitude=1.5mm, segment length=2mm}] 
          (jammer.east) -- node[below, sloped, font=\tiny, text=black, yshift=-8pt] {方向を特定} (1.2, -0.4);

        % ★修正点2：バツ印と下の文字の位置を少し下げ、全体的な被りをなくしました
        \draw[ultra thick, red] (1.1, -0.1) -- (1.7, -0.7);
        \draw[ultra thick, red] (1.1, -0.7) -- (1.7, -0.1);
        \node[font=\scriptsize\bfseries, text=red!80!black, align=center] at (2.4, -1.2) {空間的に\\遮断 (Null)};
      \end{tikzpicture}
    \end{column}
    
  \end{columns}
\end{frame}

\subsection{対策の対策}

\begin{frame}{セキュリティ学習}
  \begin{block}{セキュリティの常}
    もちろん，これらで完全に防げるわけではない．この対策にもそれにさらなる対策が講じられている．そして，これからもそのような現状が続くことが予想される．セキュリティ技術は攻撃者と対策のいたちごっこのような生存競争が常である．まだ，それを学習できる程ではないが，今後はこれらの学習を行うことを予定している．
  \end{block}
\end{frame}

\subsection{まとめ}

\begin{frame}{まとめ}
  \begin{block}{}
    様々な対策手法が存在するがそれらがすべての妨害手法をカバー可能なわけでなない．どのように対策がされていて，さらにそれを妨害者がどのようにかいくぐっているかを理解し，さらなる対策を考案する必要がある．
  \end{block}
\end{frame}

\section{現在の進捗と今後}

\subsection{利用を想定してるシミュレーション環境}

\begin{frame}{シミュレーション環境}
  妨害通信のシミュレーション環境
  \begin{itemize}
    \item サーバ環境(未定): \\
      ただし，\texttt{Dockerfile, docker-compose.yaml}は作成済みのため，おそらく即座に運用可能
    \item 使用するライブラリ(\textbf{\texttt{PyJama}}): 
      \begin{itemize}
        \item \texttt{Python3.10-slim}での利用を想定
        \item 行列演算等を用いるため，\texttt{GPU}の利用を想定
        \item githubに情報が公開(\texttt{OSS}): \texttt{https://github.com/IIP-Group/pyjama}
      \end{itemize}
  \end{itemize}
\end{frame}

\subsection{今後の学習予定}

\begin{frame}{今後の学習予定}
  \begin{itemize}
    \item 妨害の数理，仕組みについて学習
    \item 魏先生の輪講で強化学習を勉強
    \item 強化学習で用いるかつ妨害手法と対策手法の活用予想に用いるためゲーム理論を勉強

    \vspace{1em}

    \item 簡易的な妨害を行う\texttt{Python}プログラムの作成
  \end{itemize}
\end{frame}

\section{参考文献}

\begin{frame}[allowframebreaks]{参考文献}
  \begin{thebibliography}{9}
    
    % 今回のベースとなった包括的サーベイ論文（ジャミングの手口や対策の網羅的解説）
    \bibitem{Pirayesh2022}
      Hossein Pirayesh and Huacheng Zeng.
      \newblock Jamming Attacks and Anti-Jamming Strategies in Wireless Networks: A Comprehensive Survey.
      \newblock \emph{IEEE Communications Surveys \& Tutorials}, 2022.

    % OFDMシステムに対するジャミングと数学的防衛策に関する論文
    \bibitem{Yun202X} % 年度はプレプリント等に準拠
      Jaewon Yun, Joohyuk Park, and Yo-Seb Jeon.
      \newblock Anti-Jamming Modulation for OFDM Systems under Jamming Attacks.
      \newblock \emph{arXiv preprint arXiv:2510.02640}.

    % 5Gの物理層におけるセキュリティ脅威と対策に関する最新のサーベイ
    \bibitem{Harvanek2024}
      Michal Harvanek, et al.
      \newblock Survey on 5G Physical Layer Security Threats and Countermeasures.
      \newblock \emph{Sensors}, 2024.
      
    % AIやネットワークを用いた次世代の電子戦についての日本語資料
    \bibitem{Amagai2020}
      天貝 崇樹.
      \newblock 次世代の電子戦について ― 機械学習とネットワークを活用した EMS 活動 ―.
      \newblock \emph{海幹校戦略研究}, 2020.

  \end{thebibliography}
\end{frame}

\end{document}
